\documentclass[12pt]{article}
\usepackage{xr}
\usepackage[margin=1in]{geometry}
\usepackage{amsmath,amssymb,amsthm,amsfonts}
\usepackage{graphicx}
\usepackage{hyperref}
\hypersetup{colorlinks=true, linkcolor=blue, urlcolor=blue, citecolor=blue}
\theoremstyle{definition}
\newtheorem{definition}{Definition}[section]
\newtheorem{remark}{Remark}[section]
\theoremstyle{plain}
\newtheorem{theorem}{Theorem}[section]
\newtheorem{lemma}{Lemma}[section]
\newtheorem{corollary}{Corollary}[section]
\externaldocument{chapter_6}
\title{\textbf{Chapter 7: Unification of Type A and Type B measurements}}
\author{Dmytro Surdu}
\date{}

\begin{document}
\maketitle


\section{What is the notion of Type A Testable Measurement?}
\begin{definition}[Type A Testable Measurement]
A \emph{Type\,A Testable Measurement} is a quintuple
\[
  \bigl(\,
    \mathcal Q,\;
    \mathcal Z,\;
    P,\;
    T,\;
    \mathcal R
  \bigr),
\]
where
\begin{enumerate}
  \item \(\mathcal Q\) is a compact metric \emph{parameter space}.
  \item \(\mathcal Z\) is a compact metric \emph{observation space}.
  \item \(P\colon \mathcal Q \to \mathcal P(\mathcal Z)\) assigns to each \(q\in\mathcal Q\) a sampling distribution \(P_q\), and the data stream \(\{Z_k\}\) is generated i.i.d.\ from \(P_q\).
  \item \(T\colon \mathcal P(\mathcal Z)\to \Theta\subseteq \mathbb R^d\) is a Lipschitz \emph{statistic map}.
  \item \(\mathcal R\colon \Theta\to 2^{\mathcal Z}\) is Lipschitz in the Hausdorff metric, where \(\mathcal R_\theta\subseteq\mathcal Z\) is the \emph{guard‐band region} associated to statistic \(\theta\).
\end{enumerate}
The corresponding tail predicate is
\[
  P_{\rm tail}(z_{1:\infty}) 
  \;\iff\;
  \exists\,n\quad\forall\,k>n,\;
    z_k \in \mathcal R_{T(P_q)},
\]
where \(q\) is the true (unknown) parameter in \(\mathcal Q\).
\end{definition}

\section{Cross-representation of the Testable Measurement types}
\subsection{Representation of Type A Testable Measurement as a Type B}
\begin{definition}[Canonical \(\text{A}\to\text{B}\) Cast]
\label{def:canonical-A-to-B}
Let a \emph{Type A} measurement be specified by
\begin{itemize}
  \item A compact parameter space \(\mathcal Q\).
  \item An observation space \(\mathcal Z\) (compact metric).
  \item A family of sampling distributions \(\{P_q\}_{q\in\mathcal Q}\) on \(\mathcal Z\).
  \item A Lipschitz statistic map \(T:\mathcal P(\mathcal Z)\to\Theta\subseteq\mathbb R^d\).
  \item A Lipschitz guard‐band map \(\theta\mapsto\mathcal R_\theta\subseteq\mathcal Z\) in the Hausdorff metric.
\end{itemize}
We construct a \emph{Type B Testable Measurement}
\[
  \mathfrak M_{\!A}
  =\bigl(\,
     \mathcal M,\;
     \mathcal Z^{\rm in},\;
     \mathcal Z^{\rm out},\;
     E,\;
     U,\;
     \mathcal S
   \bigr)
\]
by setting
\[
  \mathcal M = \mathcal Q \times \mathcal Z,
  \qquad
  \mathcal Z^{\rm in} = \mathcal Z,
  \qquad
  \mathcal Z^{\rm out} = \mathcal Z,
\]
\[
  U\bigl((q,z),\,z'\bigr) = (\,q,z'\,),
  \qquad
  E\bigl((q,z)\bigr) = z,
\]
and defining the Certification Standard \(\mathcal S=\{\varepsilon_n\}_{n\ge1}\) by
\[
  \varepsilon_n \;=\;\frac{1}{L_T\,L_{\mathcal R}}\;\cdot\;2^{-\,\lceil \log_2 n\rceil},
\]
where \(L_T\) and \(L_{\mathcal R}\) are the Lipschitz constants of \(T\) and \(\theta\mapsto\mathcal R_\theta\).  Thus any witness of length \(n\) specifying a quantized statistic \(\widehat\theta\) induces a guard‐band \(\varepsilon_n\) that guarantees robustness of the per‐sample check
\(\;z\in\mathcal R_{\widehat\theta}\).
\end{definition}

\begin{theorem}[Type\,A \(\to\) Type\,B Cast]
\label{prop:A-to-B-cast}
Let 
\[
  (\mathcal Q,\;\mathcal Z,\;P,\;T,\;\mathcal R)
\]
be a Type\,A Testable Measurement (Definition 4.1).  Then the canonical cast
\(\mathfrak M_A\) of Definition \ref{def:canonical-A-to-B}
is a valid Type\,B Testable Measurement (Definition 3.3), and its tail predicate
agrees with the original guard‑band predicate.
\end{theorem}

\begin{proof}
Let \(\mathfrak M_A=(\mathcal M,\mathcal Z^{\rm in},\mathcal Z^{\rm out},E,U,\mathcal S)\)
be constructed as in Definition \ref{def:canonical-A-to-B}.  We verify:

\medskip\noindent\textbf{1. Spaces and maps.}
\(\mathcal M=\mathcal Q\times\mathcal Z\), \(\mathcal Z^{\rm in}=\mathcal Z^{\rm out}=\mathcal Z\)
are compact metric spaces.  The update 
\(
U((q,z),z')=(q,z')
\)
and emit
\(
E((q,z))=z
\)
maps are clearly continuous (in fact Lipschitz).

\medskip\noindent\textbf{2. Certification Standard.}
By construction \(\mathcal S=\{\varepsilon_n\}\) satisfies
the robustness requirement of Definition 3.3, since any quantized estimate
\(\widehat\theta\) of \(T(P_q)\) within \(2^{-\lceil\log_2 n\rceil}\)
induces a guard‑band \(\varepsilon_n\) guaranteeing
\(\forall z\in\mathcal Z\), \(|T(P_q)-T(P_{q'})|\) and
\(d_H(\mathcal R_{T(P_q)},\mathcal R_{T(P_{q'})})\)
remain within \(\varepsilon_n\).

\medskip\noindent\textbf{3. Tail predicate equivalence.}
Under \(\mathfrak M_A\), the belief‑state at time \(t\) is \((q,Z_t)\), so
\[
  E\bigl(U^t((q,z_0),\,z_{1:t})\bigr)=Z_t.
\]
By Chapter 6 ABET Theorem \ref{thm:basin_equivalence}, certifiability in NP is equivalent to
the existence of an attractor basin 
\(\mathcal B=\{(q,z)\!\mid\!z\in\mathcal R_{T(P_q)}\}\).
But
\[
  (q,z_{t+1:\infty})\text{ satisfies }\forall k>n\;Z_k\in\mathcal R_{T(P_q)}
  \;\Longleftrightarrow\;
  (q,z_n)\in\mathcal B\text{ and remains in }\mathcal B.
\]
This is exactly the original guard‑band predicate
\(\forall k>n,\;Z_k\in\mathcal R_{T(P_q)}\).  

Hence \(\mathfrak M_A\) is a Testable Measurement whose tail predicate
coincides with that of the given Type A measurement.
\end{proof}

\subsection{Representation of Type B Testable Measurement as a Type A}

\begin{definition}[Long‐Trajectory \(\mathrm B\!\to\!\mathrm A\) Cast]
\label{def:long-trajectory-B-to-A}
Let \(\mathfrak M=(\mathcal M,\mathcal Z^{\rm in},\mathcal Z^{\rm out},E,U,\mathcal S)\) be a Type B Testable Measurement, and let \(T_{\max}\) be its uniform reachability bound (Lemma \ref{lem:reachability-uniform}).  We define the corresponding Type A Testable Measurement
\[
  (\mathcal Q,\mathcal Z,P,T,\mathcal R)
\]
as follows:
\begin{itemize}
  \item \(\mathcal Q = \mathcal M\).
  \item \(\mathcal Z = \mathcal Z^{\rm out}\).
  \item For each \(\Theta\in\mathcal M\), let
    \[
      P_{\Theta} \;=\;\delta_{E(\Theta)},
    \]
    the Dirac measure concentrated at \(E(\Theta)\).
  \item The statistic map 
    \(
      T:\{P_{\Theta}\}\to \Theta
    \) 
    is given by 
    \(T(P_{\Theta})=\Theta\).
  \item The guard‐band map 
    \(\mathcal R:\Theta\to2^{\mathcal Z}\) 
    is defined by
    \[
      \mathcal R_{\Theta}
      = \bigl\{\,z\in\mathcal Z^{\rm out} : d\bigl(z,E(\Theta)\bigr)<\varepsilon_{T_{\max}}\bigr\},
    \]
    where \(\varepsilon_{T_{\max}}\) is the radius from the Certification Standard \(\mathcal S\) at index \(T_{\max}\).
\end{itemize}
\end{definition}

\begin{theorem}
Under the Long‐Trajectory Cast of Definition \ref{def:long-trajectory-B-to-A}, the induced tail predicate
\[
  \exists\,n\quad\forall\,k>n,\quad
    z_k\in\mathcal R_{T(P_{\Theta_0})}
\]
is equivalent to the original tail‐certification predicate “\(\Theta_0\) enters \(\mathcal B\) by time \(T_{\max}\) and thereafter \(E(\Theta_t)\) remains within the guard‐band.”  Hence the Cast is valid.
\end{theorem}

\begin{proof}
First, \(P_{\Theta}=\delta_{E(\Theta)}\) implies \(T(P_{\Theta})=\Theta\), and since \(E\) is continuous so is \(T\).  The sets \(\mathcal R_{\Theta}\) are open in \(\mathcal Z^{\rm out}\) by construction, and the map \(\Theta\mapsto\mathcal R_{\Theta}\) is Lipschitz in the Hausdorff metric because \(\mathcal S\) supplies a uniform robustness margin \(\varepsilon_{T_{\max}}\).

Next, by uniform reachability every trajectory enters \(\mathcal B\) by time \(T_{\max}\), and by definition of \(\mathcal S\) all subsequent outputs satisfy
\[
  d\bigl(Z_t,E(\Theta_{T_{\max}})\bigr)
  < \varepsilon_{T_{\max}},
  \quad t > T_{\max},
\]
so \(Z_t\in\mathcal R_{\Theta_{T_{\max}}}\).  Thus the Long‑Trajectory Cast’s guard‑band predicate holds exactly when the original Type B tail‑certificate exists.  This completes the proof.
\end{proof}


\end{document}